\section{Introduction}
\label{sec:introduction}

A forensic analyst sits before three tools and one question. The question is
simple: what URLs did the user visit on this machine? The disk image is open
in one tool, which presents a tabular view of Chrome \texttt{History} entries
extracted from the SQLite database living on the NTFS volume. The memory
dump is open in a second tool, which surfaces fragments of browser-related
strings recovered from carved pages. A third tool displays a Velociraptor
collection that pulled the live \texttt{History} file moments before the
endpoint was isolated. Three tools, three incompatible record schemas, three
different timestamp conventions, three different notions of what a ``URL
visit'' even is. The underlying artifact -- the Chromium \texttt{urls} table
and its \texttt{visits} sibling -- is identical across all three.

The concrete consequence is not abstract: the analyst must (a) write a
bespoke per-source adapter to reconcile the schemas (a multi-day engineering
task, performed under investigative time pressure, frequently re-done from
scratch in the next case), or (b) re-parse the artifact manually for each
source from raw bytes (which discards the parsing infrastructure of all
three tools), or (c) accept that evidence from one source is forensically
siloed from the others, weakening the case. None of these is acceptable; all
three occur in practice. The analyst's frustration is not a tooling
accident: it is the predictable consequence of an architectural choice that
the field has internalized without examining: \emph{evidence is classified
by what it is, not by how it is reached.}

Existing taxonomies of digital evidence are, almost without exception,
artifact-type taxonomies. Casey's textbook organizes evidence by category
(operating system, application, network, mobile)~\cite{casey2011digital}.
The CASE/UCO ontology models artifacts as typed
nodes~\cite{caseuco2018,casey2017advancing}. Carrier's data-abstraction
model layers physical, file system, and application
views~\cite{carrier2003getting,carrier2005file}. McKemmish's foundational
definition partitions forensic computing into identification, preservation,
analysis, and presentation phases against artifact
classes~\cite{mckemmish1999forensic}. Each is valuable, and each is
\emph{orthogonal} to the architectural question that determines whether the
analyst above can ever get a unified answer.

The architectural question is the address space. Disk evidence is reached by
walking a filesystem namespace from name to inode to block. Memory evidence
is reached by walking a process list, then a page table from virtual to
physical address. Log evidence is reached by seeking by timestamp or
sequence number through chunked records. Live-query evidence is reached by
dispatching a query to an endpoint and paginating result rows.
Content-addressed evidence is reached by following hashes through a Merkle
DAG. These are not five flavors of the same operation; they are five
fundamentally different ways of locating bytes -- each with distinct
dependency, trust, and durability properties. Table~\ref{tab:glance}
gives the map.

\begin{table}[h]
\centering\small
\begin{tabular}{@{}p{0.06\columnwidth}p{0.38\columnwidth}p{0.44\columnwidth}@{}}
\toprule
& \textbf{Navigation step} & \textbf{Key property} \\
\midrule
\src{P} & name $\to$ inode $\to$ block & durable by acquisition \\
\src{M} & PID $\to$ EPROCESS $\to$ VA $\to$ PA & volatile; instantaneous snapshot \\
\src{L} & timestamp $\to$ record boundary & append-only; durable by acquisition \\
\src{Q} & (endpoint, query) $\to$ result rows & produced, not stored; durable by transmission \\
\src{C} & hash $\to$ blob $\to$ Merkle DAG & immutable; intrinsically durable \\
\bottomrule
\end{tabular}
\caption{Five navigation primitives at a glance.}
\label{tab:glance}
\end{table}

We propose a taxonomy of five \emph{navigation primitives} and show that
classification by navigation primitive yields properties that artifact-type
classification cannot: parser reusability, attacker-durability profiles, and
a layered architecture whose dependency rules are mechanically enforceable.
We instantiate the framework in \emph{Issen}, a Rust platform of thirteen
crates whose layer membership is derived from the taxonomy.

\paragraph{Contributions.}
We make five claims, each tied to the section that delivers it.
\begin{enumerate}
    \item \textbf{Contribution 1 (Section~\ref{sec:taxonomy}).} The
    \emph{navigation-primitive} criterion for source classification, formalized
    as an equivalence relation $\sim$ over address-space algebras
    (Definition~\ref{def:nav-primitive}). The taxonomically consequential
    property of an evidence source is its addressing scheme, not its artifact
    category. This yields exactly five primitives under the criterion --
    \src{P}, \src{M}, \src{L}, \src{Q}, \src{C} -- defined by distinct
    address-space algebras. The three-source folk model (P/M/L) is the proper
    subset that emerges when one ignores addressing.

    \item \textbf{Contribution 2 (Section~\ref{sec:properties}, Proposition~\ref{thm:parser-indep}).}
    The \emph{Parser Independence} principle: identical bytes yield identical
    records under a fixed deterministic parser, regardless of the navigation
    primitive that retrieved the bytes. Naming it converts a folk design
    principle into an explicitly testable claim -- demonstrated empirically
    with byte-identical SHA-256 hashes across four navigation paths
    (Section~\ref{sec:case3}).

    \item \textbf{Contribution 3 (Section~\ref{sec:properties}, Propositions~\ref{prop:q-durability} and~\ref{prop:c-durability}, Corollary~\ref{cor:c-anchor-free}).}
    Formalization of \src{Q} as a first-class source type with attacker-durability
    by transmission and of \src{C} as a first-class source type with intrinsic
    attacker-durability; we prove that \src{C} is the unique source type whose
    durability requires no external trust anchor. To our knowledge, neither
    \src{Q} nor \src{C} has previously been admitted to the forensic source
    taxonomy as a peer of the disk-memory-log triad.

    \item \textbf{Contribution 4 (Section~\ref{sec:properties}, Proposition~\ref{thm:layer-rule}, and Section~\ref{sec:implementation}).}
    The \emph{Layer Dependency Rule}: a parser that imports a navigation
    component cannot be source-agnostic, because its domain is then
    $A_\pi$, not $\Bstar$. This is the contrapositive of Parser Independence.
    Issen instantiates the rule in Cargo crate boundaries enforced by the
    build system.

    \item \textbf{Contribution 5 (Sections~\ref{sec:cases}~and~\ref{sec:discussion}).}
    Four case studies, including (i) an empirical Parser Independence
    demonstration with byte-identical SHA-256 hashes across four sources,
    (ii) a counterfactual case showing concretely how a \src{P}+\src{M}+\src{L}-only
    triage tool fails to capture evidence that \src{Q} captures, and (iii)
    an explicit completeness argument framed as an open conjecture rather
    than mere absence of counterexamples.
\end{enumerate}

